% Part: proof-theory
% Chapter: natural-deduction
% Section: translation-G2i

\documentclass[../../../include/open-logic-section]{subfiles}

\begin{document}

\olfileid{pt}{nat}{tgi}

\olsection{\Log{G2i}માંથી \Log{N2i}માં રૂપાંતર}

\Log{G2i}ના સિક્વન્ટોના
પૂર્વાંગો !!{formula}sનાં
બહુગણો~$\Gamma$ છે,
જ્યારે~\Log{N2i}ના
સિક્વન્ટોના પૂર્વાંગો
એવાં ચિહ્નિત સૂત્રોના
ગણો~$\Gamma'$ છે
જેમાં દરેક ચિહ્ન માત્ર
એક વાર આવે છે.
ચિહ્નિત !!{formula}sનો
ગણ~$\Gamma'$,
બહુગણ~$\Gamma$ને
\emph{અનુરૂપ} છે એમ
ત્યારે કહીએ જ્યારે
દરેક !!{formula}~$!A$
માટે~$\Gamma'$માંના
$x : !A$ સ્વરૂપનાં
!!{element}sની સંખ્યા,
$\Gamma$માં~$!A$ની
બહુલતા (એટલે~$!A$ની
નકલોની સંખ્યા)
કરતાં વધુ ન હોય.

\begin{prop}\ollabel{prop:G2i-to-N2i}
જો $\Log{G2i} + \Cut \Proves \Gamma \Sequent \Delta$,
તો~$\Log{N2i}
\Proves \Gamma' \Sequent !C$,
જ્યાં~$\Delta$ ખાલી
હોય ત્યારે
$!C \ident \lfalse$
અને નહિતર
$\Delta = \{!A\}$ હોય,
તથા~$\Gamma'$ એ~$\Gamma$ને
અનુરૂપ હોય.
\end{prop}

\begin{proof}
  બતાવીએ કે જો~$\pi$
  એ~\Log{G2i}+\Cutમાં
  $\Gamma \Sequent \Delta$ની
  !!a{proof} હોય,
  તો~$\Gamma$ને અનુરૂપ
  એવો~$\Gamma'$ અને
  \Log{N2i}માં
  $\Gamma' \Sequent !C$ની
  !!a{proof}~$\delta$
  છે.
  \sourcecorrection{OLMD-240}{મૂળ સાબિતીના આરંભમાં અપરિભાષિત G2ci તંત્ર લખાયું છે; પ્રતિજ્ઞાપનનું તંત્ર G2i સાથે Cut છે.}
  આ માટે
  $n = \pheight{\pi}$
  પર આગમન કરીએ.

  જો~$n = 0$, તો~$\pi$
  સ્વયંસિદ્ધ છે.
  તે~$\lfalse \Sequent
  \quad$ સ્વયંસિદ્ધ હોય,
  તો~$\delta$ એ
  $x:\lfalse \Sequent \lfalse$
  સ્વયંસિદ્ધ છે,
  અર્થાત્
  $\Gamma' = \{x:\lfalse\}$
  અને~$!C \ident \lfalse$.
  જો~$\pi$ એ
  $!A \Sequent !A$
  સ્વરૂપનું સ્વયંસિદ્ધ
  હોય, તો~$\delta$
  એ~$x: !A \Sequent !A$ છે.
  
  જો~$n>0$, તો~$\pi$નો
  અંત નિયમ~$R$માં થાય છે.
  આગમન-પરિકલ્પનાથી
  \Log{N2i}માં એ
  નિયમનાં આધારવિધાનોને
  અનુરૂપ સિક્વન્ટોની
  !!{proof}s મળે છે.
  \sourcecorrection{OLMD-241}{મૂળના આગમન-વાક્યમાં N2c લખાયું છે; આ દાવો અને આખી સાબિતી N2i માટે છે.}
  જરૂરી હોય તો ફરી
  ચિહ્નિત કરીને ધારી
  શકીએ કે એ સાબિતીઓનાં
  બધાં મુક્તિ-ચિહ્નો
  જુદાં છે, અને
  કોઈ ચિહ્ન~$x$
  મુક્ત થતું હોય
  તો તેને મુક્ત કરતા
  અનુમાનની ઉપર જ
  આવે છે.
  હવે એ !!{proof}sને
  જોડીને યોગ્ય
  $\Gamma' \Sequent !C$ની
  સાબિતી કેવી રીતે
  મળે તે બતાવીએ.
  $R$ પ્રમાણે કિસ્સા
  જુદા પાડીએ.

  \begin{enumerate}
    \item જો~$R$ એ
    \RightR{\Weakening}
    હોય, તો~$\pi$નો
    અંત આ રીતે થાય છે:
    \[
    \AxiomC{}
    \RightLabel{$\pi_1$}
    \Deduce$\Gamma \fCenter$
    \RightLabel{\RightR{\Weakening}}
    \UnaryInf$\Gamma \fCenter !A$
    \DisplayProof
    \]
    $\pi_1$ પર
    આગમન-પરિકલ્પના
    લાગુ કરતાં
    $\Gamma' \Sequent \lfalse$ની
    !!a{proof} મળે છે.
    $\FalseInt$ લાગુ કરી
    $\Gamma' \Sequent !A$ની
    !!a{proof} મળે છે.
    \item જો~$R$ એ
    \LeftR{\Weakening}
    હોય, તો~$\pi$નો
    અંત આ રીતે થાય છે:
    \[
    \AxiomC{}
    \RightLabel{$\pi_1$}
    \Deduce$\Gamma \fCenter \Delta$
    \RightLabel{\LeftR{\Weakening}}
    \UnaryInf$!A, \Gamma \fCenter \Delta$
    \DisplayProof
    \]
    \sourcecorrection{OLMD-242}{મૂળના ડાબા શિથિલીકરણના વૃક્ષમાં નિયમ-લેબલ ભૂલથી જમણું શિથિલીકરણ છે; પૂર્વાંગમાં સૂત્ર ઉમેરાતું હોવાથી ડાબું શિથિલીકરણ લખ્યું છે.}
    $\pi_1$ પર
    આગમન-પરિકલ્પના
    લગાવતાં
    $\Gamma' \Sequent !C$ની
    !!a{proof} મળે છે;
    તેને~$\delta$ લઈએ.
    નોંધો કે~$\Gamma'$
    એ~$\Gamma$ને અનુરૂપ
    હોય તો~$!A,
    \Gamma$ને પણ
    અનુરૂપ છે.
    કારણ કે~$\Gamma'$માં
    $x : !A$ની સંખ્યા
    $\Gamma$માં~$!A$ની
    નકલોની સંખ્યાથી
    વધુ નથી, તેથી
    $\Gamma \cup \{!A\}$માં
    $!A$ની નકલોની
    સંખ્યાથી પણ વધુ નથી.
    \item જો~$R$ એ
    \LeftR{\Contraction}
    હોય, તો~$\pi$નો
    અંત આ રીતે થાય છે:
    \[
    \AxiomC{}
    \RightLabel{$\pi_1$}
    \Deduce$!A, !A, \Gamma \fCenter \Delta$
    \RightLabel{\LeftR{\Contraction}}
    \UnaryInf$!A, \Gamma \fCenter \Delta$
    \DisplayProof
    \]
    $\pi_1$ પર
    આગમન-પરિકલ્પના
    લગાવતાં
    $\Pi, \Gamma' \Sequent !C$ની
    !!a{proof} મળે છે,
    જ્યાં~$\Pi \subseteq \{x : !A,
    y:!A\}$.
    જો~$\Pi = \{x : !A, y:!A\}$,
    તો મળેલી~$\delta$માંના
    ચિહ્નિત
    !!{formula}s~$y:!A$
    ને~$x:!A$થી બદલો.
    $\delta$ના અંતિમ સિક્વન્ટના
    સંદર્ભમાં~$x$
    આવતો હોવાથી
    $\delta$માં કોઈ
    અનુમાન~$x:!B$
    સ્વરૂપનું
    !!a{formula} મુક્ત
    કરતું નથી એમ
    ધારી શકીએ;
    એટલે~$\delta$ના
    કોઈ સિક્વન્ટના
    ડાબે આવેલું~$x:!B$
    સક્રિય નથી.
    \sourcecorrection{OLMD-243}{મૂળમાં ચિહ્નિત N2 સૂત્રોનું નામબદલણ G2i ઉપસાબિતી $\pi_1$માં કરવાનું કહેવાયું છે; આગમનથી મળેલી N2 સાબિતી $\delta$માં તે કરવું જોઈએ.}
    તેથી પરિણામ
    \Log{N2i}માં
    હજી પણ !!a{proof} છે.
    \item જો~$R$ એ
    $\RightR{\lif}$ હોય,
    તો~$\pi$નો અંત
    આ રીતે થાય છે:
    \[
    \AxiomC{}
    \RightLabel{$\pi_1$}
    \Deduce$!A, \Gamma \fCenter !B$
    \RightLabel{\RightR{\lif}}
    \UnaryInf$\Gamma \fCenter !A \lif !B$
    \DisplayProof
    \]
    આગમન-પરિકલ્પનાથી
    $x:
    !A, \Gamma' \Sequent !B$
    અથવા~$\Gamma' \Sequent !B$ની
    !!a{proof}~$\delta_1$
    છે, જ્યાં~$\Gamma'$
    એ~$\Gamma$ને
    અનુરૂપ છે.
    $\delta_1$ પછી
    \Intro{\lif}
    લગાવી મળતી
    સાબિતીને~$\delta$
    લઈ શકાય.
    \sourcecorrection{OLMD-244}{મૂળમાં $\delta_1$ એ $\delta$ પછી અનુસરણ-પરિચય લગાવી મળે છે એમ ઉલટું લખાયું છે; આગમનથી $\delta_1$ મળે છે અને તેની નીચે નિયમ લગાવી $\delta$ બને છે.}
    \item જો~$R$ એ
    \LeftR{\lif} હોય,
    તો~$\pi$નો
    અંત આ રીતે થાય છે:
    \[
    \AxiomC{}
    \RightLabel{$\pi_1$}
    \Deduce$\Gamma_1 \fCenter !A$
    \AxiomC{}
    \RightLabel{$\pi_2$}
    \Deduce$!B, \Gamma_2 \fCenter \Delta$
    \RightLabel{\LeftR{\lif}}
    \BinaryInf$!A \lif !B, \Gamma_1, \Gamma_2 \fCenter \Delta$
    \DisplayProof
    \]
    આગમન-પરિકલ્પનાથી
    $\Gamma_1' \Sequent !A$ની
    !!{proof}~$\delta_1$
    અને~$x: !B, \Gamma_2' \Sequent !C$
    અથવા~$\Gamma_2' \Sequent !C$ની
    સાબિતી~$\delta_2$
    મળે છે.
    $\delta_2$ બીજા
    સ્વરૂપની હોય,
    તો~$\delta_2$ને જ~$\delta$
    લઈ શકાય.
    નહિતર~$\delta_1$માંનાં
    બધા સૂત્રો તથા
    મુક્તિ-ચિહ્નો ફરી
    ચિહ્નિત કરીને
    $\delta_1$ અને~$\delta_2$
    વચ્ચે કોઈ ચિહ્ન
    સામાન્ય ન રહે
    તે સુનિશ્ચિત કરીએ.
    \sourcecorrection{OLMD-245}{મૂળમાં આ N2 ચિહ્નબદલી G2i ઉપસાબિતી $\pi_1$માં કરવાનું કહેવાયું છે; ચિહ્નિત સૂત્રો અને મુક્તિ-ચિહ્નો ધરાવતી અનુવાદિત સાબિતી $\delta_1$ છે.}
    ત્યારે~$\Gamma_1' \cup
    \Gamma_2'$ એ
    $\Gamma_1 \cup \Gamma_2$
    ને અનુરૂપ છે.
    $x:!B$ એ~$\delta_2$ના
    અંતિમ સિક્વન્ટમાં
    હોવાથી~$\delta_2$માં
    $x$ મુક્તિ-ચિહ્નવાળા
    કોઈ અનુમાનમાં
    $x:!B$ સક્રિય
    હોઈ શકે નહીં.
    $\delta_2$માંનું
    દરેક સ્વયંસિદ્ધ
    $x:!B \Sequent !B$
    નીચેની
    !!{proof}~$\delta_3$થી
    બદલો:
    \[
    \Axiom$x: !A \lif !B \fCenter !A \lif !B$
    \AxiomC{}
    \RightLabel{$\delta_1$}
    \Deduce$\Gamma_1' \fCenter !A$
    \RightLabel{\Elim{\lif}}
    \BinaryInf$x: !A \lif !B, \Gamma_1' \fCenter !B$
    \DisplayProof
    \]
    \sourcecorrection{OLMD-246}{મૂળના $\delta_3$ અનુસરણ-નિવારણ વૃક્ષના નિષ્કર્ષમાં $\Gamma_1'$ છૂટે છે; બીજું આધારવિધાન એ સંદર્ભ ધરાવે છે અને નિયમ તેને નિષ્કર્ષમાં રાખે છે.}
    અને~$\delta_2$માં
    $x:!B$ના દરેક
    સ્થાનને
    $x:!A \lif !B$થી
    બદલો.
    પરિણામ
    $\Subst{\delta_2}{\delta_3}{x:!B}$
    એ
    $x:!A \lif !B, \Gamma_1', \Gamma_2' \Sequent !C$ની
    !!a{proof} છે.
    \item જો~$R$ એ
    \RightR{\land} હોય,
    તો~$\pi$નો
    અંત આ રીતે થાય છે:
    \[
    \AxiomC{}
    \RightLabel{$\pi_1$}
    \Deduce$\Gamma_1 \fCenter !A$
    \AxiomC{}
    \RightLabel{$\pi_2$}
    \Deduce$\Gamma_2 \fCenter !B$
    \RightLabel{\RightR{\land}}
    \BinaryInf$\Gamma_1, \Gamma_2 \fCenter !A \land !B$
    \DisplayProof
    \]
    આગમન-પરિકલ્પનાથી
    $\Gamma_1' \Sequent !A$ની
    સાબિતી~$\delta_1$
    અને~$\Gamma_2' \Sequent !B$ની
    સાબિતી~$\delta_2$
    મળે છે.
    $\delta_2$માંનાં
    બધા સૂત્રો તથા
    મુક્તિ-ચિહ્નો ફરી
    ચિહ્નિત કરીને
    $\delta_1$ અને~$\delta_2$
    વચ્ચે કોઈ ચિહ્ન
    સામાન્ય ન રહે
    તે સુનિશ્ચિત કરીએ.
    \sourcecorrection{OLMD-247}{મૂળમાં બીજા આગમન-નિષ્કર્ષનું સૂત્ર $!A$ લખાયું છે, પણ બીજા આધારવિધાનમાં $!B$ છે; વળી N2 ચિહ્નબદલી $\pi_2$ને બદલે $\delta_2$માં થવી જોઈએ.}
    ત્યારે~$\Gamma_1' \cup \Gamma_2'$
    એ~$\Gamma_1 \cup
    \Gamma_2$ને અનુરૂપ છે.
    $\delta_1$ તથા~$\delta_2$ના
    અંતિમ સિક્વન્ટો પર
    $\Intro{\land}$
    લગાવી
    !!a{proof}~$\delta$
    મળે છે.
    \item જો~$R$ એ
    $\LeftR{\land}$ હોય,
    તો, દાખલા તરીકે,
    $\pi$નો અંત આ રીતે થાય છે:
    \[
    \AxiomC{}
    \RightLabel{$\pi_1$}
    \Deduce$!A, \Gamma \fCenter \Delta$
    \RightLabel{\LeftR{\land}}
    \UnaryInf$!A \land !B, \Gamma \fCenter \Delta$
    \DisplayProof
    \]
    આગમન-પરિકલ્પનાથી
    $x:
    !A, \Gamma' \Sequent !C$
    અથવા~$\Gamma' \Sequent !C$ની
    !!a{proof}~$\delta_1$
    છે, જ્યાં~$\Gamma'$
    એ~$\Gamma$ને
    અનુરૂપ છે.
    બીજા કિસ્સામાં
    $\delta$ને~$\delta_1$
    લઈ શકાય.
    \sourcecorrection{OLMD-248}{મૂળના બીજા કિસ્સામાં $\delta_1$ને $\delta$ કહેવાનું વાક્ય દિશા ઉલટી બતાવે છે; આગમનથી મળેલી $\delta_1$ જ માંગેલી $\delta$ તરીકે લેવાય છે.}
    પહેલા કિસ્સામાં
    $x:!A$ અંતિમ
    સિક્વન્ટમાં હોવાથી
    $\delta_1$માં~$x$
    મુક્તિ-ચિહ્નવાળા
    કોઈ અનુમાનમાં
    $x:!A$ સક્રિય
    નથી. દરેક
    સ્વયંસિદ્ધ
    $x:!A
    \Sequent !A$ને
    આ~!!{proof}~$\delta_2$થી
    બદલો:
    \[
    \Axiom$x: !A \land !B \fCenter !A \land !B$
    \RightLabel{\Elim{\land}}
    \UnaryInf$x: !A \land !B \fCenter !A$
    \DisplayProof
    \]
    અને~$x:!A$ના
    દરેક સ્થાનને
    $x:!A \land !B$થી
    બદલો; એટલે
    $\Subst{\delta_1}{\delta_2}{x:!A}$
    લો. આ
    $x:!A
    \land !B, \Gamma' \Sequent !C$ની
    !!a{proof} છે.
    \item જો~$R$ એ
    \Cut હોય, તો~$\pi$નો
    અંત આ રીતે થાય છે:
    \[
    \AxiomC{}
    \RightLabel{$\pi_1$}
    \Deduce$\Gamma_1 \fCenter !A$
    \AxiomC{}
    \RightLabel{$\pi_2$}
    \Deduce$!A, \Gamma_2 \fCenter \Delta$
    \RightLabel{\Cut}
    \BinaryInf$\Gamma_1, \Gamma_2 \fCenter \Delta$
    \DisplayProof
    \]
    આગમન-પરિકલ્પનાથી
    $\Gamma_1' \Sequent !A$ની
    સાબિતી~$\delta_1$
    અને
    $x:!A, \Gamma_2' \Sequent !C$
    અથવા~$\Gamma_2' \Sequent !C$ની
    સાબિતી~$\delta_2$
    મળે છે.
    બીજા કિસ્સામાં
    $\delta$ને~$\delta_2$
    લો. નહિતર
    $\delta_2$માંનાં
    દરેક સ્વયંસિદ્ધ
    $x:!A \Sequent !A$
    ને~$\delta_1$થી
    બદલો. આમ
    $\Gamma_1',
    \Gamma_2' \Sequent !C$ની
    !!{proof}
    $\Subst{\delta_2}{\delta_1}{x:!A}$
    એ~$\delta$ બને છે.
  \end{enumerate}
\end{proof}

\end{document}
